% ---- score2 standard preamble (from: python3 score2_tools.py preamble) ----
\documentclass[10pt]{extarticle}
\usepackage[a4paper,margin=1.5cm]{geometry}
\usepackage[T1]{fontenc}
\usepackage{amsmath,amssymb}
\usepackage{graphicx}
\usepackage[export]{adjustbox}
\usepackage{multicol}
\usepackage{enumitem}
\usepackage{booktabs,array,tabularx}
\setlength{\columnsep}{0.7cm}
\setlength{\parindent}{0pt}
\setlength{\parskip}{3pt}
\pagestyle{empty}
\makeatletter
% figure registry (figdefs.tex is generated by the extract step)
\newcommand{\qfigdef}[4]{\@namedef{qf@file@#1}{#2}\@namedef{qf@w@#1}{#3}\@namedef{qf@wide@#1}{#4}}
\newcommand{\qf@check}[1]{\@ifundefined{qf@file@#1}{\PackageError{score2}{Unknown figure id #1}{}}{}}
\newcommand{\qf@img}[4][]{\includegraphics[width=#2\dimexpr\@nameuse{qf@w@#4}pt\relax,max width=#3,max height=0.45\textheight,#1]{\@nameuse{qf@file@#4}}}
% Full-width questions: finalize wraps every question that contains a wide figure in \qfullbegin...\qfullend,
% so the WHOLE question (stem, figures, options) is set across the page in reading order, then the two columns
% resume. The question's head - its number, stem and first figure (\qfullsplit marks the end of the head) -
% is kept on one page, so a number is never left at a page foot with its figure overleaf; the rest may break
% across pages normally, so no big gaps are left. Without a \qfullsplit, the whole question is kept together
% if it is under 45% of a page. Never write these yourself.
\newif\ifqf@full
\newif\ifqf@boxed
\newsavebox{\qf@box}
\newcommand{\qf@open}{\global\qf@fulltrue\par\global\qf@boxedtrue
  \setbox\qf@box\vbox\bgroup\hsize\textwidth\linewidth\textwidth}
\newcommand{\qf@emit}[1]{\par\egroup\global\qf@boxedfalse
  \ifdim\dimexpr\ht\qf@box+\dp\qf@box\relax<#1\textheight\noindent\usebox{\qf@box}\par\else\unvbox\qf@box\fi}
\newcommand{\qfullsplit}{\ifqf@boxed\qf@emit{0.85}\fi}
\newcommand{\qf@close}{\ifqf@boxed\qf@emit{0.45}\fi\global\qf@fullfalse}
\newcommand{\qfullbegin}{\end{multicols}\qf@open}
\newcommand{\qfullend}{\qf@close\begin{multicols}{2}}
\newcommand{\qfullstart}{\qf@open}
\newcommand{\qfulldone}{\qf@close}
% \qfig[scale]{id}: block figure, centred, directly where it is written (size set by the tool).
% Fallback if a wide figure is not inside a full-width question: break out of the columns in place - no floats.
\newcommand{\qfig}[2][1]{\qf@check{#2}%
  \ifqf@full\qf@blk{#1}{#2}\else
  \ifnum\@nameuse{qf@wide@#2}=1 \expandafter\@firstoftwo\else\expandafter\@secondoftwo\fi
  {\end{multicols}\begin{center}\qf@img{#1}{\textwidth}{#2}\end{center}\begin{multicols}{2}}%
  {\qf@blk{#1}{#2}}\fi}
\newcommand{\qf@blk}[2]{\par\smallskip{\centering\qf@img{#1}{\linewidth}{#2}\par}\smallskip}
% \qfigi[scale]{id}: inline figure (inside option lists / tables / minipages); never breaks columns.
\newcommand{\qfigi}[2][1]{\qf@check{#2}\qf@img[valign=c]{#1}{\linewidth}{#2}}
\makeatother
% structure macros - lint reads these to check completeness and figure placement
\newcommand{\papersec}[1]{\par\bigskip{\centering\Large\bfseries PAPER-#1\par}\medskip}
\newcommand{\partsec}[2]{\par\medskip\noindent{\large\bfseries PART-#1: #2}\par\smallskip}
\newcommand{\sect}[1]{\par\smallskip\noindent\textbf{#1}\par}
\newcommand{\qn}[1]{\par\smallskip\noindent\textbf{#1.}~}
\newcommand{\paraq}[2]{\par\smallskip\noindent\textbf{Paragraph for Questions #1 and #2}\par}
\newenvironment{qoptions}{\begin{itemize}[leftmargin=1.8em,itemsep=1pt,topsep=2pt,label={}]}{\end{itemize}}
\newcommand{\opt}[2]{\item[(#1)] #2}
% figure registry (generated by score2_tools.py finalize)
\qfigdef{p1_phy_s1_q01_1}{Score-I - Test 22 (04 Dec 2022)-img/p1_phy_s1_q01_1.png}{100.1}{0}
\qfigdef{p1_phy_s1_q02_1}{Score-I - Test 22 (04 Dec 2022)-img/p1_phy_s1_q02_1.png}{137.4}{0}
\qfigdef{p1_phy_s1_q04_1}{Score-I - Test 22 (04 Dec 2022)-img/p1_phy_s1_q04_1.png}{159.6}{0}
\qfigdef{p1_phy_s1_q06_1}{Score-I - Test 22 (04 Dec 2022)-img/p1_phy_s1_q06_1.png}{100.7}{0}
\qfigdef{p1_phy_s1_q07_2}{Score-I - Test 22 (04 Dec 2022)-img/p1_phy_s1_q07_2.png}{92.3}{0}
\qfigdef{p1_phy_s1_q07_1}{Score-I - Test 22 (04 Dec 2022)-img/p1_phy_s1_q07_1.png}{116.5}{0}
\qfigdef{p1_phy_s1_q08_1}{Score-I - Test 22 (04 Dec 2022)-img/p1_phy_s1_q08_1.png}{161.1}{0}
\qfigdef{p1_phy_s2_para9_10_1}{Score-I - Test 22 (04 Dec 2022)-img/p1_phy_s2_para9_10_1.png}{114.4}{0}
\qfigdef{p1_phy_s2_para9_10_2}{Score-I - Test 22 (04 Dec 2022)-img/p1_phy_s2_para9_10_2.png}{121.2}{0}
\qfigdef{p1_phy_s2_para9_10_3}{Score-I - Test 22 (04 Dec 2022)-img/p1_phy_s2_para9_10_3.png}{120.1}{0}
\qfigdef{p1_phy_s2_para9_10_4}{Score-I - Test 22 (04 Dec 2022)-img/p1_phy_s2_para9_10_4.png}{133.7}{0}
\qfigdef{p1_phy_s2_para11_12_1}{Score-I - Test 22 (04 Dec 2022)-img/p1_phy_s2_para11_12_1.png}{150.6}{0}
\qfigdef{p1_phy_s2_q12_1}{Score-I - Test 22 (04 Dec 2022)-img/p1_phy_s2_q12_1.png}{128.5}{0}
\qfigdef{p1_phy_s3_q01_1}{Score-I - Test 22 (04 Dec 2022)-img/p1_phy_s3_q01_1.png}{101.9}{0}
\qfigdef{p1_phy_s3_q02_1}{Score-I - Test 22 (04 Dec 2022)-img/p1_phy_s3_q02_1.png}{120.7}{0}
\qfigdef{p1_phy_s3_q03_1}{Score-I - Test 22 (04 Dec 2022)-img/p1_phy_s3_q03_1.png}{151.4}{0}
\qfigdef{p1_phy_s3_q05_1}{Score-I - Test 22 (04 Dec 2022)-img/p1_phy_s3_q05_1.png}{69.4}{0}
\qfigdef{p1_phy_s3_q05_2}{Score-I - Test 22 (04 Dec 2022)-img/p1_phy_s3_q05_2.png}{88.7}{0}
\qfigdef{p1_phy_s3_q06_1}{Score-I - Test 22 (04 Dec 2022)-img/p1_phy_s3_q06_1.png}{142.0}{0}
\qfigdef{p1_chem_s1_q01_1}{Score-I - Test 22 (04 Dec 2022)-img/p1_chem_s1_q01_1.png}{142.8}{0}
\qfigdef{p1_chem_s1_q02_A}{Score-I - Test 22 (04 Dec 2022)-img/p1_chem_s1_q02_A.png}{25.2}{0}
\qfigdef{p1_chem_s1_q02_B}{Score-I - Test 22 (04 Dec 2022)-img/p1_chem_s1_q02_B.png}{48.7}{0}
\qfigdef{p1_chem_s1_q02_C}{Score-I - Test 22 (04 Dec 2022)-img/p1_chem_s1_q02_C.png}{40.9}{0}
\qfigdef{p1_chem_s1_q02_D}{Score-I - Test 22 (04 Dec 2022)-img/p1_chem_s1_q02_D.png}{26.2}{0}
\qfigdef{p1_chem_s1_q04_A}{Score-I - Test 22 (04 Dec 2022)-img/p1_chem_s1_q04_A.png}{64.9}{0}
\qfigdef{p1_chem_s1_q04_B}{Score-I - Test 22 (04 Dec 2022)-img/p1_chem_s1_q04_B.png}{65.1}{0}
\qfigdef{p1_chem_s1_q04_C}{Score-I - Test 22 (04 Dec 2022)-img/p1_chem_s1_q04_C.png}{63.4}{0}
\qfigdef{p1_chem_s1_q04_D}{Score-I - Test 22 (04 Dec 2022)-img/p1_chem_s1_q04_D.png}{46.6}{0}
\qfigdef{p1_chem_s1_q05_1}{Score-I - Test 22 (04 Dec 2022)-img/p1_chem_s1_q05_1.png}{123.7}{0}
\qfigdef{p1_chem_s1_q06_1}{Score-I - Test 22 (04 Dec 2022)-img/p1_chem_s1_q06_1.png}{84.6}{0}
\qfigdef{p1_chem_s1_q06_2}{Score-I - Test 22 (04 Dec 2022)-img/p1_chem_s1_q06_2.png}{88.2}{0}
\qfigdef{p1_chem_s1_q07_A}{Score-I - Test 22 (04 Dec 2022)-img/p1_chem_s1_q07_A.png}{58.2}{0}
\qfigdef{p1_chem_s1_q07_B}{Score-I - Test 22 (04 Dec 2022)-img/p1_chem_s1_q07_B.png}{56.8}{0}
\qfigdef{p1_chem_s1_q07_C}{Score-I - Test 22 (04 Dec 2022)-img/p1_chem_s1_q07_C.png}{49.5}{0}
\qfigdef{p1_chem_s1_q07_D}{Score-I - Test 22 (04 Dec 2022)-img/p1_chem_s1_q07_D.png}{32.6}{0}
\qfigdef{p1_chem_s1_q08_1}{Score-I - Test 22 (04 Dec 2022)-img/p1_chem_s1_q08_1.png}{42.4}{0}
\qfigdef{p1_chem_s1_q08_2}{Score-I - Test 22 (04 Dec 2022)-img/p1_chem_s1_q08_2.png}{48.7}{0}
\qfigdef{p1_chem_s1_q08_C}{Score-I - Test 22 (04 Dec 2022)-img/p1_chem_s1_q08_C.png}{59.3}{0}
\qfigdef{p1_chem_s1_q08_3}{Score-I - Test 22 (04 Dec 2022)-img/p1_chem_s1_q08_3.png}{56.3}{0}
\qfigdef{p1_chem_s2_para9_10_1}{Score-I - Test 22 (04 Dec 2022)-img/p1_chem_s2_para9_10_1.png}{59.0}{0}
\qfigdef{p1_chem_s2_para9_10_2}{Score-I - Test 22 (04 Dec 2022)-img/p1_chem_s2_para9_10_2.png}{54.1}{0}
\qfigdef{p1_chem_s2_para9_10_3}{Score-I - Test 22 (04 Dec 2022)-img/p1_chem_s2_para9_10_3.png}{44.1}{0}
\qfigdef{p1_chem_s2_para9_10_4}{Score-I - Test 22 (04 Dec 2022)-img/p1_chem_s2_para9_10_4.png}{58.3}{0}
\qfigdef{p1_chem_s2_para11_12_1}{Score-I - Test 22 (04 Dec 2022)-img/p1_chem_s2_para11_12_1.png}{42.4}{0}
\qfigdef{p1_chem_s2_para11_12_2}{Score-I - Test 22 (04 Dec 2022)-img/p1_chem_s2_para11_12_2.png}{27.7}{0}
\qfigdef{p1_chem_s2_para11_12_3}{Score-I - Test 22 (04 Dec 2022)-img/p1_chem_s2_para11_12_3.png}{38.0}{0}
\qfigdef{p1_chem_s2_para11_12_4}{Score-I - Test 22 (04 Dec 2022)-img/p1_chem_s2_para11_12_4.png}{60.8}{0}
\qfigdef{p1_chem_s3_q01_1}{Score-I - Test 22 (04 Dec 2022)-img/p1_chem_s3_q01_1.png}{223.0}{0}
\qfigdef{p1_chem_s3_q02_1}{Score-I - Test 22 (04 Dec 2022)-img/p1_chem_s3_q02_1.png}{85.4}{0}
\qfigdef{p1_chem_s3_q03_1}{Score-I - Test 22 (04 Dec 2022)-img/p1_chem_s3_q03_1.png}{222.4}{0}
\qfigdef{p1_chem_s3_q04_1}{Score-I - Test 22 (04 Dec 2022)-img/p1_chem_s3_q04_1.png}{205.2}{0}
\qfigdef{p1_chem_s3_q05_1}{Score-I - Test 22 (04 Dec 2022)-img/p1_chem_s3_q05_1.png}{229.5}{0}
\qfigdef{p1_chem_s3_q06_1}{Score-I - Test 22 (04 Dec 2022)-img/p1_chem_s3_q06_1.png}{70.5}{0}
% ---- body skeleton ----
% \begin{document}
% \begin{center}{\Large\bfseries JEE (Main+Advanced) Score-II --- Test NN (DD Mon YYYY)}\end{center}
% \begin{multicols}{2}
%   \papersec{1}   (only for two-paper tests)
%   \partsec{1}{PHYSICS} \sect{Section-I (i): Single Correct} \qn{1} ... \qfig{p1_phy_q01_1} ...
% \end{multicols}
% \end{document}

\begin{document}
\begin{center}{\Large\bfseries JEE (Main+Advanced) Score-I --- Test 22 (04 Dec 2022)}\end{center}
\raggedcolumns
\setlength{\tabcolsep}{3pt}
\begin{multicols}{2}
\papersec{1}
\partsec{1}{PHYSICS}
\sect{SECTION-I (i)}
\qn{1} Two points masses $A$ and $B$ each of 1 kg joined by 1 m rod (light) translates on horizontal surface as shown in figure. A particle of mass 1 kg at rest on horizontal surface sticks to ball as it collides with ball $A$. then :
\qfig{p1_phy_s1_q01_1}
\begin{qoptions}
\opt{A}{Linear speed of ball $A$ just after collision is 2 m/s}
\opt{B}{Linear speed of ball $B$ just after collision is 4 m/s}
\opt{C}{Velocity of centre of mass of system is 8/3 m/s}
\opt{D}{Angular speed of system, about centre of mass after collision is 2 rad/s.}
\end{qoptions}

\qn{2} As shown in the figure, $ABC$ is a triangular plate of mass $M$ and side length $L$ each. Its mass distribution is non-uniform but its centre of mass is at the centroid. Moment of inertia of this plate about an axis passing through its centroid and perpendicular to its plane is given by $\dfrac{ML^2}{24}$. End $C$ of the plate is connected to a massless and inextensible string of length $L$ as shown. At the other end of the thread a point particle of mass $M$ is connected. At $t=0$ the point mass is given initial velocity of $v_0$ perpendicular to the length of the string. Choose the correct option(s).
\qfig{p1_phy_s1_q02_1}
\begin{qoptions}
\opt{A}{The radius of curvature of the path of the point mass just after projection is $3L$.}
\opt{B}{The radius of curvature of the path of the point mass just after projection is $4L$.}
\opt{C}{The tension in the string just after projection is $\dfrac{Mv_0^2}{4L}$}
\opt{D}{Net acceleration of end $B$ and $C$ of plate just after the projection will be same in magnitude.}
\end{qoptions}

\qn{3} A point object starts moving in a circle of radius 1 mm with constant angular acceleration of 1 rad/s$^2$ in a plane perpendicular to the principal axis at a distance of 20 cm from the pole of a concave mirror of focal length 25 cm. At time $t=2$ s after the motion starts
\begin{qoptions}
\opt{A}{The ratio of acceleration of image and the object is $\sqrt{29}$.}
\opt{B}{The ratio of acceleration of image and the object is 5.}
\opt{C}{Magnitude of relative velocity of object with respect to image is 8 mm/s}
\opt{D}{Magnitude of relative velocity of object with respect to image is 12 mm/s}
\end{qoptions}

\qn{4} A uniform circular ring of mass $m=5$ kg and radius $R=30$ cm is hinged at the point $A$ of a stationary surface with the line joining point $A$ to the centre $C$ of the ring is initially vertical. The ring is also in contact with the Plank of mass $m=5$ kg at point $B$ which is placed on a smooth horizontal surface. Initially points $A$ and $B$ are almost touching and the line $AC$ is vertical and the system is released from rest. Then the velocity of the Plank when the line $AC$ makes angle $30^\circ$ with the horizontal as shown will be (there is no friction between the ring and the plank and $g=10$ m/s$^2$)
\qfig{p1_phy_s1_q04_1}
\begin{qoptions}
\opt{A}{1 m/s}
\opt{B}{2 m/s}
\opt{C}{3 m/s}
\opt{D}{4 m/s}
\end{qoptions}

\qn{5} A ray of light moving along vector $\left(5\sqrt2\,\hat i-5\hat j-5\hat k\right)$ undergoes refraction at an interface of two media, which is $y$-$z$ plane. The refractive index for $x\le0$ is $\sqrt3$ while for $x\ge0$ it is $\sqrt2$. Then the unit vector along the refracted ray is :
\begin{qoptions}
\opt{A}{$\dfrac{\sqrt3\,\hat i-\sqrt2\,\hat j-\sqrt2\,\hat k}{\sqrt8}$}
\opt{B}{$\dfrac{\sqrt3\,\hat i-\sqrt3\,\hat j-\sqrt2\,\hat k}{\sqrt8}$}
\opt{C}{$\dfrac{\sqrt2\,\hat i-\sqrt3\,\hat j-\sqrt3\,\hat k}{\sqrt8}$}
\opt{D}{$\dfrac{\sqrt2\,\hat i+\sqrt3\,\hat j-\sqrt3\,\hat k}{\sqrt8}$}
\end{qoptions}

\qn{6} A particle of mass $m=1$ kg is rigidly attached to the circumference at point $A$ of a uniform circular ring of mass $m=1$ kg and radius $R=20$ cm placed on a rough inclined plane of inclination $\theta=37^\circ$ as shown in the figure. Then choose the correct option(s) just after the ring is released from rest. (there is no slipping between the ring and inclined plane and $g=10$ m/s$^2$)
\qfig{p1_phy_s1_q06_1}
\begin{qoptions}
\opt{A}{The frictional force acting on the ring due to inclined plane is 4N.}
\opt{B}{The frictional force acting on the ring due to inclined plane is 2N.}
\opt{C}{The normal force acting on the ring due to inclined plane is 22N.}
\opt{D}{The normal force acting on the ring due to inclined plane is 11N.}
\end{qoptions}

\qn{7} A glass prism is immersed in a hypothetical liquid. The graph shows the variation of refractive index of glass and liquid with the wavelength $\lambda$ of light. When a ray of white light is incident parallel to the base of the prism, then select the correct option
\begin{center}\qfigi{p1_phy_s1_q07_2}\hfill\qfigi{p1_phy_s1_q07_1}\end{center}
\begin{qoptions}
\opt{A}{Yellow ray travels without deviation.}
\opt{B}{Blue ray is deviated towards vertex.}
\opt{C}{Red ray is deviated towards base.}
\opt{D}{There is no dispersion.}
\end{qoptions}

\qn{8} Two rods each or mass $m$ and length $\dfrac{\ell}{4}$ are hinged on a larger square plate of same mass and edge length $2\ell$. Both rods can freely rotate about hinge points $P$ and $Q$ in the plane of square plate. Choose the correct option(s).
\qfig{p1_phy_s1_q08_1}
\begin{qoptions}
\opt{A}{The moment of inertia of the system about $x$-axis is $\dfrac{65}{192}m\ell^2$.}
\opt{B}{The moment of inertia of the system about $y$-axis is $\dfrac{161}{192}m\ell^2$.}
\opt{C}{The moment of inertia of the system about $z$-axis is $\dfrac{226}{192}m\ell^2$.}
\opt{D}{The moment of inertia of the system about $z$-axis is $\dfrac{224}{192}m\ell^2$.}
\end{qoptions}

\end{multicols}
\noindent\begin{minipage}{\textwidth}
\sect{SECTION-I (ii)}
\paraq{9}{10}
The figures show an isosceles right angled prism of ($\mu=1.6$) and there is a point source situated either inside or outside of the prism also we have two observers $P$ \& $Q$. Opaque sheet having a small hole lies outside the prism. List-I represents location of observers and List-II represents the nature \& location of images.
\par\smallskip
\begin{center}\small
\renewcommand{\arraystretch}{1.1}
\begin{tabular}{|c|p{10.0cm}|c|p{4.2cm}|}
\hline
 & \multicolumn{1}{c|}{List-I} & & \multicolumn{1}{c|}{List-II}\\
\hline
(I) & \qfigi[0.85]{p1_phy_s2_para9_10_1}\hspace{0.8em}\parbox[c]{4.6cm}{A cavity is made inside the prism in the shape of equi-convex lens of radius of curvature 20 cm and the cavity is filled with a transparent liquid of refractive index $\mu=2.4$. Source is outside the prism as shown} & (P) & Only P can see the image\\
\hline
(II) & \qfigi[0.85]{p1_phy_s2_para9_10_2}\hspace{0.8em}\parbox[c]{4.6cm}{Source is outside the prism as shown and there no any cavity inside the prism.} & (Q) & Only Q can see the image\\
\hline
\end{tabular}
\end{center}
\end{minipage}
\noindent\begin{minipage}{\textwidth}
\begin{center}\small
\renewcommand{\arraystretch}{1.1}
\begin{tabular}{|c|p{10.0cm}|c|p{4.2cm}|}
\hline
(III) & \qfigi[0.85]{p1_phy_s2_para9_10_3}\hspace{0.8em}\parbox[c]{4.6cm}{A cavity is made inside the prism in the shape of equi-convex lens of radius of curvature 20 cm and the cavity is filled with a transparent liquid of refractive index $\mu=2.4$. Source is inside the prism as shown} & (R) & Image is virtual located at $55/8$ cm above the surface $BC$\\
\hline
(IV) & \qfigi[0.85]{p1_phy_s2_para9_10_4}\hspace{0.8em}\parbox[c]{4.6cm}{Source is inside the prism as shown and there is no cavity inside the prism.} & (S) & Image is real\\
\hline
 & & (T) & No observer can see the image\\
\hline
 & & (U) & Both (P) \& (Q) can see the image\\
\hline
\end{tabular}
\end{center}
\end{minipage}
\begin{multicols}{2}
\qn{9} Which of following is correct combination if medium outside the prism have refractive index $\mu=1$.
\begin{qoptions}
\opt{A}{I -- R}
\opt{B}{II -- R}
\opt{C}{I -- P}
\opt{D}{III -- P}
\end{qoptions}

\qn{10} Which of following is correct combination if outside medium is $\mu=1.6$.
\begin{qoptions}
\opt{A}{I -- S}
\opt{B}{II -- R}
\opt{C}{III -- S}
\opt{D}{IV -- Q}
\end{qoptions}

\paraq{11}{12}
\qfig{p1_phy_s2_para11_12_1}
The figure shows two identical rollable objects held at rest and the spring shown is unstretched and unattached to any object. The complete surface is sufficiently rough to sustain pure rolling all the time. Mass of each object is $m$ and radius $r$ is very small compared to the height $h$. Spring constant of spring is $k$. The object kept at height is released so that it goes on to compress the spring. Possible rollable objects are given in List-I. Some values are given in List-II.
\begin{center}\small
\renewcommand{\arraystretch}{1.25}
\begin{tabular}{|c|c|c|c|}
\hline
 & List-I & & List-II\\
\hline
(I) & Solid sphere & (P) & $\sqrt{\dfrac{3m}{4k}}$\\
\hline
(II) & Hollow sphere & (Q) & $\sqrt{\dfrac{3m}{2k}}$\\
\hline
(III) & Ring & (R) & $\sqrt{\dfrac{m}{k}}$\\
\hline
(IV) & Disc & (S) & $\sqrt{\dfrac{5m}{3k}}$\\
\hline
 & & (T) & $\sqrt{\dfrac{5m}{6k}}$\\
\hline
 & & (U) & $\sqrt{\dfrac{7m}{10k}}$\\
\hline
\end{tabular}
\end{center}
\qn{11} The correct match for the ratio of maximum compression in spring to maximum speed of the object $A$ is
\begin{qoptions}
\opt{A}{I-P, II-R, III-U, IV-S}
\opt{B}{I-R, II-U, III-S, IV-T}
\opt{C}{I-U, II-T, III-R, IV-P}
\opt{D}{I-S, II-R, III-Q, IV-U}
\end{qoptions}

\qn{12} If a smooth fixed vertical wall exists to the right of the lower object touching it. The correct match for the ratio $\dfrac{\text{maximum compression in spring}}{\sqrt2\,(\text{maximum speed of the objective }A)}$ is
\qfig{p1_phy_s2_q12_1}
\begin{qoptions}
\opt{A}{I--T, II--P, III--S, IV--Q}
\opt{B}{I--Q, II--S, III--P, IV--T}
\opt{C}{I--R, II--Q, III--P, IV--S}
\opt{D}{I--U, II--T, III--R, IV--P}
\end{qoptions}

\sect{SECTION-II}
\qn{1} $O$ is a small object placed at a distance $D$ from eye of an observer. A concave lens of focal length $f$ ($f<D$) is placed near to eye and image of the object is viewed. Now the lens is moved towards object $O$ away from the eye. Find the distance (in m) of lens from object when appeared size of image is smallest ($D=6$ m, $f=20$ cm)
\qfig{p1_phy_s3_q01_1}

\qn{2} Bahubali tossed a rod of length $L=10$ m from horizontal position at $t=0$ such that its centre receives a velocity $v=\pi\sqrt{gL}$ in vertically upward direction and some angular velocity about its centre. The rod lands back to its point of projection in horizontal position. During its course of flight its angular velocity remains constant and it made one complete rotation. Calculate the time (in sec) after which end $B$ is at its highest vertical position given that rod rotates in vertical plane. Take solution of equation $\cos\theta=2\theta$ to be $\theta=0.45$ and $\sin(0.45)=0.43$ if needed.
\qfig{p1_phy_s3_q02_1}

\qn{3} The final image $I$ of the object $O$ shown in the arrangement is formed at point 20 cm below a thin equi-convex lens, which is at a depth of 65 cm from principal axis. Calculate the radius of curvature in cm of lens kept at $A$.
\qfig{p1_phy_s3_q03_1}

\qn{4} In a setup of displacement method experiment, distance between the screen and a light source is 100 cm and the lens used has a small aperture. By moving the lens between the source and the screen, sharp images are obtained on the screen for two different positions of the lens separated by a distance $x$ cm. The ratio of sizes of these two images is 9 : 1. The focal length of the lens used is $f$ cm. The value $x$ is :

\qn{5} A thin convex lens of focal length 50 cm is cut into three parts $A$, $B$ and $C$ by the planes parallel to its optic axis. The thickness of the middle layer $C$ is 9 mm. The middle layer is now removed and the two parts $A$ and $B$ are stuck together. Then the part $C$ is also placed in contact of this lens symmetrically as shown in the figure. Now, a point object $O$ is placed at a distance 100 cm left on the optic axis of the part $C$. Find the distance (in mm) between the two images formed by the parts $A$ and $B$.
\begin{center}\qfigi{p1_phy_s3_q05_1}\hfill\qfigi{p1_phy_s3_q05_2}\end{center}

\qn{6} A luminous point object $O$ is placed in front of spherical glass of radius $R=20$ cm. The final real image is formed after refraction through the spherical surface, followed by reflection at the concave mirror of radius $R=400$ cm and second refraction through the spherical surface. Find the distance (in cm) of this final image from the pole $P_1$ of the spherical surface.
\qfig{p1_phy_s3_q06_1}

\partsec{2}{CHEMISTRY}
\sect{SECTION-I (i)}
\qn{1} \qfigi{p1_chem_s1_q01_1}\\
Select the correct statement(s) regarding the product $A$ \& $B$.
\begin{qoptions}
\opt{A}{B is V}
\opt{B}{A is IV}
\opt{C}{A and B both are III}
\opt{D}{A is II and B is I}
\end{qoptions}

\qn{2} In which of the following species, the positive or negative charge is involved in resonance.
\begin{qoptions}
\opt{A}{\qfigi{p1_chem_s1_q02_A}}
\opt{B}{\qfigi{p1_chem_s1_q02_B}}
\opt{C}{\qfigi{p1_chem_s1_q02_C}}
\opt{D}{\qfigi{p1_chem_s1_q02_D}}
\end{qoptions}

\qn{3} $\underset{\text{(C)}}{\text{C}_8\text{H}_{14}}$ $\xrightarrow[\text{(1 eqvt)}]{\text{H}_2/\text{Pd}}$ $\underset{\text{(D)}}{\text{C}_8\text{H}_{16}}$ $\xrightarrow[\text{Zn}]{\text{O}_3}$ 2--butanone (only product formed)\\
Select the correct statement(s) regarding above reaction sequence.
\begin{qoptions}
\opt{A}{Compound `C' may be an alkyne or a conjugated diene.}
\opt{B}{Compound `D' shows markonikov addition with HBr \& anti markonikov addition with HBr/ROOR.}
\opt{C}{2-butanone has no chain isomers.}
\opt{D}{Compound `D' shows geometrical isomerism but doesn't show optical isomerism.}
\end{qoptions}

\qn{4} One mole of the following compound(s) can consume more than 4 moles of HI.
\begin{qoptions}
\opt{A}{\qfigi{p1_chem_s1_q04_A}}
\opt{B}{\qfigi{p1_chem_s1_q04_B}}
\opt{C}{\qfigi{p1_chem_s1_q04_C}}
\opt{D}{\qfigi{p1_chem_s1_q04_D}}
\end{qoptions}

\qn{5} Select the correct statement(s) for the given reaction is/are.
\qfig{p1_chem_s1_q05_1}
\begin{qoptions}
\opt{A}{`B' is an epoxide}
\opt{B}{`A' is cis-alkene}
\opt{C}{`B' does not have chiral centers.}
\opt{D}{`A' is produced via anti-elimination.}
\end{qoptions}

\qn{6} \qfigi{p1_chem_s1_q06_1}\hspace{1.5em}\qfigi{p1_chem_s1_q06_2}\par
\begin{qoptions}
\opt{A}{(i) and (iv) are enantiomers}
\opt{B}{(i) and (iii) are diastereomers}
\opt{C}{(ii) and (iii) are meso compounds}
\opt{D}{(ii) and (iii) are diastereomers.}
\end{qoptions}

\qn{7} Which of the following species gives CO$_2$ with NaHCO$_3$?
\begin{qoptions}
\opt{A}{\qfigi{p1_chem_s1_q07_A}}
\opt{B}{\qfigi{p1_chem_s1_q07_B}}
\opt{C}{\qfigi{p1_chem_s1_q07_C}}
\opt{D}{\qfigi{p1_chem_s1_q07_D}}
\end{qoptions}

\qn{8} Select the incorrect statements
\begin{qoptions}
\opt{A}{Rate of hydrolysis of \qfigi{p1_chem_s1_q08_1} is slower than that of \qfigi{p1_chem_s1_q08_2}, via SN$^1$ mechanism.}
\opt{B}{Rate of solvolysis of (CH$_3$)$_3$CBr is faster in H$_2$O than in C$_2$H$_5$OH.}
\opt{C}{\qfigi{p1_chem_s1_q08_C} reacts faster than \qfigi{p1_chem_s1_q08_3} in both SN$^2$ reactions.}
\opt{D}{CH$_3$O--CH=CH--CH$_2$Br reacts slower than H$_2$C=C(OCH$_3$)--CH$_2$Br in SN$^1$ solvolysis reaction}
\end{qoptions}

\sect{SECTION-I (ii)}
\paraq{9}{10}
In the following, List-I contains some organic compounds, and List-II contains properties.
\par\smallskip
\begin{center}\small
\begin{tabular}{|c|c|c|p{4.2cm}|}
\hline
 & List-I & & List-II\\
\hline
(I) & \qfigi{p1_chem_s2_para9_10_1}\rule{0pt}{26pt} & (P) & Aromatic compound\\
\hline
(II) & \qfigi{p1_chem_s2_para9_10_2}\rule{0pt}{26pt} & (Q) & More acidic than C$_2$H$_5$OH\\
\hline
(III) & \qfigi{p1_chem_s2_para9_10_3}\rule{0pt}{26pt} & (R) & Undergoes addition reaction with electrophiles\\
\hline
(IV) & \qfigi{p1_chem_s2_para9_10_4}\rule{0pt}{26pt} & (S) & Undergoes substitution reaction with odd electron species\\
\hline
 & & (T) & Produces gas with Grignard reagent.\\
\hline
 & & (U) & Soluble in NaOH\\
\hline
\end{tabular}
\end{center}
\qn{9} Which of the following options has the correct combination considering List-I and List-II.
\begin{qoptions}
\opt{A}{(II) -- P, Q, R, S, T, U}
\opt{B}{(III) -- P, R}
\opt{C}{(II) -- P, Q, R, S, T}
\opt{D}{(III) -- R, S, T}
\end{qoptions}

\qn{10} Which of the following options has the correct combination considering List-I and List-II.
\begin{qoptions}
\opt{A}{(I) -- P, S, U}
\opt{B}{(IV) -- P, T}
\opt{C}{(I) -- P, S, T}
\opt{D}{(IV) -- S}
\end{qoptions}

\paraq{11}{12}
List-I contains four isomers of C$_7$H$_{10}$O and List-II contains reactions \& properties.
\par\smallskip
\begin{center}\small
\begin{tabular}{|c|c|c|p{4.4cm}|}
\hline
 & List-I & & List-II\\
\hline
(I) & \qfigi{p1_chem_s2_para11_12_1}\rule{0pt}{22pt} & (P) & Produces H$_2$ on reaction with metallic Na\\
\hline
(II) & \qfigi{p1_chem_s2_para11_12_2}\rule{0pt}{22pt} & (Q) & Undergoes oxidation to aldehyde or ketone by PCC\\
\hline
(III) & \qfigi{p1_chem_s2_para11_12_3}\rule{0pt}{22pt} & (R) & Cloudiness appears in Lucas Test at room temperature\\
\hline
(IV) & \qfigi{p1_chem_s2_para11_12_4}\rule{0pt}{22pt} & (S) & Produces CH$_4$ with CH$_3$MgBr\\
\hline
 & & (T) & forms H--bond with H$_2$O\\
\hline
 & & (U) & Gives red color in Victor Meyer's test.\\
\hline
\end{tabular}
\end{center}
\qn{11} Which of the following options has the correct combination considering List-I and List-II.
\begin{qoptions}
\opt{A}{(I) -- P, Q, S, T, U}
\opt{B}{(I) -- P, Q, R, S, T}
\opt{C}{(IV) -- S, T, U}
\opt{D}{(IV) -- R, T, U}
\end{qoptions}

\qn{12} Which of the following options has the correct combination considering List-I and List-II.
\begin{qoptions}
\opt{A}{(II) -- P, Q, R, T}
\opt{B}{(II) -- P, R, T, U}
\opt{C}{(III) -- P, Q, R, S, T}
\opt{D}{(III) -- Q, R, S, T, U}
\end{qoptions}

\sect{SECTION-II}
\qn{1} Calculate the number of stereoisomers for the following compound :
\qfig{p1_chem_s3_q01_1}

\qn{2} How many of the following species are electrophilic in nature.\\
\qfigi{p1_chem_s3_q02_1}, SO$_3$, CH$_3$--C$\equiv$O$^+$, Ph--$\overset{+}{\text N}\equiv$N, Cl$_2$, MnO$_4^-$

\qn{3} How many of the following molecules show geometrical as well as optical isomerism.\\
(i) pent--2,3--diene, \quad (ii) octa--2,3,6--trine, \quad (iii) 1,3--dimethyl cyclohexane
\qfig{p1_chem_s3_q03_1}

\qn{4} How many of the following compounds are soluble in aqueous NaOH.
\qfig{p1_chem_s3_q04_1}

\qn{5} How many of the following compounds are ortho -- para directing in electrophilic aromatic substitution.
\qfig{p1_chem_s3_q05_1}

\qn{6} Find the total number of pairs of enantiomers and pairs of diastereomers obtained upon complete ozonolysis of the following compound
\qfig{p1_chem_s3_q06_1}

\partsec{3}{MATHEMATICS}
\sect{SECTION-I (i)}
\qn{1} Let $L$ be line containing by the planes $P_1: x+3y=2$ and $P_2: 3x-y+z=0$. If $C$ lies on $L$ and $A\equiv(5,-1,-5)$, $B\equiv(1,1,-2)$, then
\begin{qoptions}
\opt{A}{shortest distance between lines $L$ and $AB$ is $\sqrt{\dfrac{22}{57}}$}
\opt{B}{shortest distance between lines $L$ and $AB$ is $\sqrt{\dfrac{88}{57}}$}
\opt{C}{minimum area of $\Delta ABC$ is $2\sqrt{\dfrac{319}{114}}$ sq. units}
\opt{D}{minimum area of $\Delta ABC$ is $\sqrt{\dfrac{319}{114}}$ sq. units}
\end{qoptions}

\qn{2} Consider,\\
$f(x)=15\sin^{-1}(\sin x)-{}$\\
$\Big[3+\big[1+[1+\{x\}]\{x\}\big]\{x\}\Big]\,x$\\
where $[k]$ and $\{k\}$ denotes greatest integer function and fractional part function of $k$ respectively. Identify which of the following statement(s) is(are) correct?
\begin{qoptions}
\opt{A}{Number of solutions of the equation $f(x)=0$ in $[-\pi,\pi]$ is 3}
\opt{B}{Number of solutions of the equation $f(x)=0$ in $[-3\pi,3\pi]$ is 5}
\opt{C}{Sum of the solutions of the equation $f(x)=0$ in $[-2\pi,3\pi]$ is $\dfrac{5\pi}{2}$}
\opt{D}{Sum of the solutions of the equation $f(x)=0$ in $[-\pi,5\pi]$ is $\dfrac{5\pi}{2}$}
\end{qoptions}

\qn{3} If $A=\{\theta:2\cos^2\theta+\sin\theta\le2\}$ and $B=\{\theta:\sin^{-1}(\cos\theta)\le\cos^{-1}(\cos\theta)\}$ where $\theta\in[0,2\pi]$, then
\begin{qoptions}
\opt{A}{number of integers in $A\cap B=4$}
\opt{B}{$\left[\dfrac\pi4,\dfrac{5\pi}{6}\right]\subseteq A\cap B$}
\opt{C}{$\left[\pi,\dfrac{3\pi}{2}\right]\subseteq A\cap B$}
\opt{D}{$A\cup B=\left[\dfrac\pi6,2\pi\right]$}
\end{qoptions}

\qn{4} Let $\vec a,\vec b,\vec c$ be three vectors such that $(|\vec a|-1)^2+\left(\left|\vec b\right|-2\right)^2+(|\vec c|-3)^2\le0$ and $\left|\vec a-2\vec b+\vec c\right|^2+\left|\vec a+\vec b-2\vec c\right|^2+\left|-2\vec a+\vec b+\vec c\right|^2=126$, then
\begin{qoptions}
\opt{A}{$\vec a,\vec b,\vec c$ constitute a triangle}
\opt{B}{$\cos(\hat a\wedge\vec b)+\cos(\vec b\wedge\vec c)+\cos(\vec c\wedge\vec a)=1$}
\opt{C}{$\left(\vec a\times\vec b\right)\times\vec c=3\left(2\vec a-\vec b\right)$}
\opt{D}{$\left|\vec a+2\vec b+4\vec c\right|=7$}
\end{qoptions}

\qn{5} Identify which of the following statement is(are) correct on the set $A=\{a,b,c,d\}$
\begin{qoptions}
\opt{A}{$R=\{(a,a),(b,b),(c,c),(d,d),(a,b),(b,a)\}$ is an equivalence relation on $A$.}
\opt{B}{$R=\{(a,a),(b,b),(c,c),(d,d),(a,b)\}$ is transitive relation on $A$.}
\opt{C}{Number of symmetric relations on $A$ is equal to $2^{10}$.}
\opt{D}{Number of reflexive relations on $A$ is equal to $2^{14}$.}
\end{qoptions}

\qn{6} If $A(1,1,1)$ and $C(9,10,12)$ are the two vertices of a triangle $ABC$ and equation of bisector $\angle ABC$ is $\dfrac{x-1}{4}=\dfrac{z-6}{-2}$ \& $y=4$ then
\begin{qoptions}
\opt{A}{Equation of side $BC$ is $\dfrac{x-5}{4}=\dfrac{y-7}{3}=\dfrac{z-9}{3}$}
\opt{B}{Equation of side $BC$ is $\dfrac{x-2}{3}=\dfrac{y-3}{4}=\dfrac{z-4}{5}$}
\opt{C}{Co-ordinates of vertex $B$ is $(1,5,2)$}
\opt{D}{Co-ordinates of vertex $B$ is $(1,4,6)$}
\end{qoptions}

\qn{7} A relation $U$ is defined as $U=\{(x,y):x^y=y^x;\ x,y\in N\}$ relation $U$ is
\begin{qoptions}
\opt{A}{Symmetric}
\opt{B}{Transitive}
\opt{C}{Reflexive}
\opt{D}{Void relation}
\end{qoptions}

\qn{8} Let $f:R^+\to R$, $f(\lambda)=1$, where $\lambda\in R^+$ and $f(x)\,f(y)+f\!\left(\dfrac\lambda x\right)f\!\left(\dfrac\lambda y\right)=2f(xy)\ \forall\ x,y\in R^+$, then
\begin{qoptions}
\opt{A}{range of $f(x)$ contains only two integers}
\opt{B}{$f(x)$ is periodic function}
\opt{C}{$\displaystyle\sum_{r=-50}^{50}f\!\left(\dfrac{51-r}{101-r}\right)=101$, $r\in I$}
\opt{D}{$f(x)$ is bounded function}
\end{qoptions}

\sect{SECTION-I (ii)}
\paraq{9}{10}
In $\Delta ABC$, $\begin{cases}\overrightarrow{AB}=\hat i+\sqrt\alpha\,\hat j+\alpha\hat k\\ \overrightarrow{AC}=\hat i-\sqrt\alpha\,\hat j+\alpha\hat k\end{cases}$ and area of $\Delta ABC=\Delta$.\\
There are some expressions are given in List-I whose values are given in List-II below.
\par\smallskip
\begin{center}\small
\renewcommand{\arraystretch}{1.2}
\begin{tabular}{|c|p{4.2cm}|c|c|}
\hline
 & \centering List-I & & List-II\\
\hline
(I) & greatest value of $\angle BAC$ & (P) & 1\\
\hline
(II) & The length of side $\overrightarrow{BC}$, when $\angle BAC$ is maximum & (Q) & 2\\
\hline
(III) & The least value of $|\Delta/\alpha|$ is & (R) & 3\\
\hline
(IV) & The least value of $|\Delta/\alpha|$ occurs when $\alpha$ is & (S) & $\sqrt2$\\
\hline
 & & (T) & $\tan^{-1}\!\left(2\sqrt2\right)$\\
\hline
 & & (U) & $\sin^{-1}\!\left(\dfrac13\right)$\\
\hline
\end{tabular}
\end{center}
\qn{9} Which of the following is only the \textbf{CORRECT} combination?
\begin{qoptions}
\opt{A}{I $\to$ U}
\opt{B}{II $\to$ Q}
\opt{C}{III $\to$ R}
\opt{D}{IV $\to$ Q}
\end{qoptions}

\qn{10} Which of the following is only the \textbf{CORRECT} combination?
\begin{qoptions}
\opt{A}{I $\to$ T}
\opt{B}{II $\to$ R}
\opt{C}{III $\to$ P}
\opt{D}{IV $\to$ Q}
\end{qoptions}

\paraq{11}{12}
Consider the function $y=\sin^{-1}x$, $y=\tan^{-1}x$, $y=\cos^{-1}\!\left(\dfrac{1-x^2}{1+x^2}\right)$ and $y=\sin^{-1}\!\left(\dfrac{2x}{1+x^2}\right)$. The possible integral values of `$a$' for which the graph $y=|x-a|$ will cut the graph of given functions is mentioned in List-II.
\par\smallskip
\begin{center}\small
\renewcommand{\arraystretch}{1.3}
\begin{tabular}{|c|c|c|c|}
\hline
 & List-I & & List-II\\
\hline
(I) & $y=\sin^{-1}x$ & (P) & $a=-2$\\
\hline
(II) & $y=\tan^{-1}x$ & (Q) & $a=-1$\\
\hline
(III) & $y=\cos^{-1}\!\left(\dfrac{1-x^2}{1+x^2}\right)$ & (R) & $a=0$\\
\hline
(IV) & $y=\sin^{-1}\!\left(\dfrac{2x}{1+x^2}\right)$ & (S) & $a=1$\\
\hline
 & & (T) & $a=2$\\
\hline
 & & (U) & $a=3$\\
\hline
\end{tabular}
\end{center}
\qn{11} Which of the following is \textbf{CORRECT} combination?
\begin{qoptions}
\opt{A}{(I) $\to$ P,Q,R,S}
\opt{B}{(II) $\to$ P,Q}
\opt{C}{(I) $\to$ R,S,T}
\opt{D}{(II) $\to$ P,Q,R}
\end{qoptions}

\qn{12} Which of the following is \textbf{INCORRECT} combination?
\begin{qoptions}
\opt{A}{(III) $\to$ R,S,T,U}
\opt{B}{(IV) $\to$ R,S,T,U}
\opt{C}{(III) $\to$ P,Q,R}
\opt{D}{(II) $\to$ Q,R,S,T}
\end{qoptions}

\sect{SECTION-II}
\qn{1} A rectangle $ABCD$ of dimensions $\sqrt{85}$ and $2\sqrt{85}$ is folded along the diagonal $BD$ such that planes $ABD$ and $CBD$ are perpendicular to each other. Let the position of the vertex $A$ remains unchanged and $C_1$ is the new position of $C$. Then the distance of $C_1$ from $A$ is equal to

\qn{2} The number of solution(s) to equation $2^x=x^2$, $x\in R$ is $k$, then fundamental period of the function $f(x)=kx-[kx]+[|\sin(\cos x)|]$ where $[t]$ denotes greatest integer function less than or equal to $t$, is

\qn{3} If the range of the function $f(x)=(\sin^{-1}x)^3+(\cos^{-1}x)^3+\left(\sin^{-1}\!\left(4x^4-4x^2+1\right)\right)^3$ is $[a\pi^3,b\pi^3]$, then $\left(\dfrac{8a}{b}\right)$ is equal to

\qn{4} The maximum volume of parallelepiped having vectors along co-terminal edges as $\vec a\times\vec b$, $2\vec b\times\vec c$, $3\vec c\times\vec a$, where $\vec a=(\sin^2\alpha)\hat i+2\hat j-3\hat k$; $\vec b=\hat i+2(\sin^2\alpha)\hat j-2\hat k$; $\vec c=\hat i+\hat j-(\sin^2\alpha)\hat k$, $\alpha\in R$ is

\qn{5} If `$a$' is number of value(s) of `$x$' satisfying $3^x+4^x+5^x+\ldots\ldots+n^x-(n+1)^x=0$ where $n$ is a natural number and $3<n<100$ and $f$ is number of integers in the range of the function, $y=\sin^{-1}(\sin^{-1}x)$. Then minimum value of the expression $\left|e^ax^2+[f-e]|x|+e\right|$; $x\in R$ is\\
\textbf{Note:} $[t]$ denotes greatest integer function less than or equal to $t$ and $e$ is napier's constant.

\qn{6} Total number of students appearing for campus placements from IIT X, IIT Y and IIT Z are 290, 340 and 360 respectively.\\
Company A can only hire from IIT X and IIT Y.\\
Company B can only hire from IIT Y and IIT Z.\\
Company C can only hire from IIT Z and IIT X.\\
If total hiring done by company A, B and C are 280, 300 and 185 respectively.\\
IIT X and IIT Y secured exactly 90\% and 80\% placement respectively \& `$N$' is the number of students remains unplaced, then the value of $\dfrac{N}{10}$ is
\end{multicols}
\end{document}
